\documentclass[11pt]{article}

\usepackage[T1]{fontenc}
\usepackage{microtype}
\usepackage[a4paper,margin=28mm]{geometry}
\usepackage{amsmath,amssymb,amsthm,mathtools}
\usepackage{hyperref}

\hypersetup{
  hidelinks,
  pdftitle={A Rigorous Reduction for the Two-Player 3n plus or minus 1 Game},
  pdfauthor={Grisha Pochuev}
}

\newtheorem{theorem}{Theorem}[section]
\newtheorem{proposition}[theorem]{Proposition}
\newtheorem{lemma}[theorem]{Lemma}
\newtheorem{corollary}[theorem]{Corollary}
\theoremstyle{definition}
\newtheorem{definition}[theorem]{Definition}
\newtheorem{remark}[theorem]{Remark}

\newcommand{\odd}{\operatorname{odd}}
\newcommand{\Win}{\mathsf{WIN}}
\newcommand{\Loss}{\mathsf{LOSS}}
\newcommand{\Draw}{\mathsf{DRAW}}
\newcommand{\N}{\mathbb N}
\newcommand{\vTwo}{v_2}
\newcommand{\xor}{\mathbin{\mathtt{xor}}}
\newcommand{\proved}{\textsc{proved}}
\newcommand{\verified}{\textsc{verified by computation}}
\newcommand{\conditional}{\textsc{conditional theorem}}
\newcommand{\openstatus}{\textsc{open target, not claimed}}

\title{A Rigorous Reduction for the Two-Player $3n\boldsymbol{\pm}1$ Game\\
\large Verified Lemmas, a Conditional No-Draw Theorem, and the Remaining Refinement Problem}
\author{Grisha Pochuev}
\date{Adversarial-audit edition, 12 August 2026}

\begin{document}
\maketitle

\begin{center}
\fbox{%
\begin{minipage}{0.91\textwidth}
\textbf{Scope notice.}
This manuscript does \emph{not} claim an unconditional solution of the game.
It gives complete proofs of the exact binary normal form, the outcome lemmas,
the finite-token rank, and the abstract no-draw criterion.  The application of
that criterion to every legal game continuation still requires a concrete
semantic refinement.  One unresolved token-provenance defect in the current
assembly, and a second defect in the earlier v6.3 router, are identified in
Section~\ref{sec:audit}.  Consequently, the assertion that
optimal play always terminates is marked \openstatus, not \proved.  Within
this stated reduction-and-audit scope, every theorem asserted by the
manuscript has a proof and the final adversarial audit reports no red or
orange defect.
\end{minipage}}
\end{center}

\begin{abstract}
From an odd integer $n>1$, a player chooses $3n+1$ or $3n-1$ and removes all
powers of two; reaching $1$ wins.  Arbitrary play may cycle, so the problem is
to exclude game-theoretic draws, not merely to find a decreasing move.  We
derive a self-contained conjugacy to the binary game
\[
  q\longmapsto A(q)=\left\lceil\frac{3q}{2}\right\rceil,
  \qquad
  q\longmapsto B(q)=R(A(q)),
\]
where $R$ deletes the maximal alternating binary suffix and $B(q)<q$ for
$q>0$.  We prove the exact finite-outcome recursion without imposing a finite
search boundary, give a Gray-code description of $R$, and prove a
well-founded multiset rank for finite proof witnesses.  We then state and
prove an explicit semantic-refinement theorem: a productive lift of every
draw into a certified lexicographically decreasing macro system rules out all
draws.  The repository's routing certificate checks only the abstract macro
system.  It does not construct the required lift and successor map.  A
hostile audit exposes an uninstalled-witness error in the current high-return
rank.  It also records an outcome-orientation error in the superseded v6.3
loss-sibling router, so that the history of the argument is not silently
rewritten.  These are gaps in proposed global bridges, not counterexamples to
the game theorem.  The result of
this paper is therefore a rigorous reduction with a precise open obligation.
The shortest marked return is sharpened further: its apparent valuation-one
subcase is empty, so that row always replaces an already retained losing token
by one of its actual children.
\end{abstract}

\tableofcontents

\section{Claim ledger and scope}

Every theorem or lemma in this paper followed by a proof is labelled
\proved.  Finite experiments are labelled \verified.  The global target is
not inferred from either kind of finite computation.

\begin{center}
\begin{tabular}{p{0.55\textwidth}p{0.33\textwidth}}
\hline
Claim & Status \\
\hline
Exact $m$-space and $q$-space normal forms & \proved \\
$B(q)<q$ and $B(B(A(q)))<q$ for $q>0$ & \proved \\
Finite $\Win/\Loss$ recursion and the draw-child lemma & \proved \\
Gray-code formula for the alternating-suffix deletion & \proved \\
Finite-token replacement strictly lowers the rank & \proved \\
Complete semantic macro refinement implies no draws & \conditional \\
The checked JSON inventory supplies that refinement & \openstatus \\
Optimal play terminates from every odd start & \openstatus \\
Finite identities and regression suites listed in Section~\ref{sec:computation}
  & \verified \\
Fixed-tail loss-sibling and source-zero $j=1$ local repairs & \proved \\
Shortest marked return ($D=1$) strictly pays its retained token & \proved \\
\hline
\end{tabular}
\end{center}

The distinction in the last four rows is the central audit boundary.  A
certificate may correctly prove that a declared transition system is
well-founded while failing to prove that every relevant game situation has a
legal, outcome-compatible transition in that system.

\section{The game and its infinite-graph semantics}

For a positive even integer $x$, let $\vTwo(x)$ be the exponent of $2$ in
$x$, and put
\[
  \odd(x)=\frac{x}{2^{\vTwo(x)}}.
\]
The positions are the positive odd integers.  Position $1$ is terminal and is
losing for the player to move.  For odd $n>1$, the two children are
\[
  T_-(n)=\odd(3n-1),\qquad T_+(n)=\odd(3n+1).
\]
The graph is infinite and may contain cycles; for example
$5\to7\to5$ is legal.  Finite-game backward induction may therefore not be
used with an artificial boundary.

\subsection{Finite proof trees}

\begin{definition}[Finite outcome layers]
Let $L_0=\{1\}$ and $W_0=\varnothing$.  Inductively, for $h\ge0$, define
\begin{align*}
 W_{h+1}&=W_h\cup
 \{n>1:\text{ at least one child of $n$ belongs to $L_h$}\},\\
 L_{h+1}&=L_h\cup
 \{n>1:\text{ both children of $n$ belong to $W_h$}\}.
\end{align*}
Put $W=\bigcup_hW_h$, $L=\bigcup_hL_h$, and
$D=\{1,3,5,\ldots\}\setminus(W\cup L)$.
Membership in $W$ or $L$ is witnessed by a finite proof tree.  The least layer
of membership is its proof height.
\end{definition}

\begin{lemma}[Disjointness; \proved]\label{lem:disjoint}
$W\cap L=\varnothing$.
\end{lemma}

\begin{proof}
Suppose otherwise and choose a position $x$ with a pair of $W$- and
$L$-derivations for which the sum of the two derivation heights is minimal.
The terminal position is never in $W$, so $x>1$.  The $W$-derivation chooses
a child $y$ having an $L$-derivation of smaller height.  The $L$-derivation
requires every child, including $y$, to have a $W$-derivation of smaller
height.  Thus $y$ has both labels with a smaller sum of heights, a
contradiction.
\end{proof}

\begin{definition}[Outcomes]
We call the positions in $W$, $L$, and $D$ respectively $\Win$, $\Loss$, and
$\Draw$.  The word ``draw'' means failure of both finite winning
certificates.  It is not the same as ``not reached by a bounded search''.
\end{definition}

\begin{proposition}[Exact recursion; \proved]\label{prop:recursion}
For every nonterminal position $x$:
\begin{enumerate}
\item[(i)] $x$ is $\Win$ if and only if at least one child is $\Loss$;
\item[(ii)] $x$ is $\Loss$ if and only if both children are $\Win$;
\item[(iii)] if $x$ is $\Draw$, then at least one child is $\Draw$ and no child is
      $\Loss$.
\end{enumerate}
\end{proposition}

\begin{proof}
The forward implications in (i) and (ii) follow from the last rule in a
minimal finite derivation.  Conversely, a finite proof for a losing child may
be extended by one node to a winning proof for its parent; two finite winning
proofs may be joined below a losing parent.  This proves (i) and (ii).
If a draw had a losing child, (i) would make it winning.  If neither child
were a draw, both would be resolved; since neither is losing, both would be
winning, and (ii) would make the parent losing.  This proves (iii).
\end{proof}

\begin{proposition}[Game-theoretic meaning; \proved]\label{prop:strategy}
At a $\Win$ position, the player to move has a strategy reaching $1$ after
finitely many moves.  At a $\Loss$ position, every first move gives the
opponent such a strategy.  At a $\Draw$ position neither player can force a
win.  Hence, if $D=\varnothing$, optimal play terminates at $1$ from every
starting position.
\end{proposition}

\begin{proof}
Induct on the finite proof height.  At a winning node, move to the certified
losing child.  At a losing node, either move leads to a winning node of lower
proof height in its attached finite proof tree.  Along play consistent with
the relevant winning strategy the heights decrease through a finite tree, so
the terminal node is reached.

Now start at a draw.  Proposition~\ref{prop:recursion}(iii) supplies a draw
child and excludes losing children.  Against any proposed winning strategy,
the opponent preserves a draw whenever the proposed strategy leaves a draw;
if it instead moves to a winning child (winning for the new player to move),
the opponent switches to that child's finite winning strategy.  Thus the
proposed strategy cannot force a win.  The same argument applies to either
player.  The final assertion follows from Lemma~\ref{lem:disjoint} and the
definition of $D$.
\end{proof}

\section{Exact binary normal form}

Write an odd position as $n=2m+1$, $m\ge0$, and define
\[
  F(m)=\left\lceil\frac{3m}{2}\right\rceil.
\]
For $m>0$, let $R(m)$ be the integer obtained by deleting the maximal
alternating suffix of the ordinary binary expansion of $m$.  If the whole
word is alternating, the result is $0$; set $R(0)=0$.  For example,
\[
  R(1001_2)=10_2,\qquad R(10101_2)=0.
\]

The following recurrence is an exact definition suitable for induction:
\begin{equation}\label{eq:Rrec}
\begin{aligned}
R(2a)&=\begin{cases}a,&a\text{ even},\\R(a),&a\text{ odd},\end{cases}\\
R(2a+1)&=\begin{cases}R(a),&a\text{ even},\\a,&a\text{ odd}.
\end{cases}
\end{aligned}
\end{equation}
Indeed, the low bit extends the alternating suffix exactly when it differs
from the preceding bit.

\begin{theorem}[$m$-space normal form; \proved]\label{thm:mnormal}
For $n=2m+1>1$, the two children, expressed again as
$n'=2m'+1$, have coordinates
\[
  \boxed{m'=F(m)\quad\text{and}\quad m'=F(R(m)).}
\]
The unordered pair is meant; the sign producing each member depends on the
parity of $m$.
\end{theorem}

\begin{proof}
Since
\[
  3n-1=2(3m+1),\qquad 3n+1=2(3m+2),
\]
exactly one of $3m+1$ and $3m+2$ is odd.  If $m=2a$, the odd expression is
$6a+1$ and its new coordinate is $3a=F(m)$.  If $m=2a+1$, the odd expression
is $6a+5$ and its new coordinate is $3a+2=F(m)$.

It remains to identify the branch whose inner expression is even.  We use
strong induction on $m$.  For $m=0$ the only relevant position is terminal.
If $m=2a$, the other odd child is $\odd(3a+1)$.  When $a$ is even, $3a+1$ is
odd, its coordinate is $3a/2=F(a)$, and~\eqref{eq:Rrec} gives $R(m)=a$.
When $a$ is odd, $3a+1$ is the even branch occurring in the smaller input
$a$; the induction hypothesis gives coordinate $F(R(a))$, and
$R(m)=R(a)$.  If $m=2a+1$, the other odd child is $\odd(3a+2)$.  When $a$ is
odd it has coordinate $(3a+1)/2=F(a)$ and $R(m)=a$; when $a$ is even, the
case $a=0$ gives coordinate $0=F(R(0))$ directly, while for $a>0$ the
induction hypothesis at $a$ gives $F(R(a))$ and $R(m)=R(a)$.  These four cases
prove the formula.
\end{proof}

\subsection{Conjugation}

The image of $F$ is precisely the set of nonnegative integers not congruent
to $1\pmod3$, and
\begin{equation}\label{eq:Finverse}
 F^{-1}(y)=
 \begin{cases}
  2y/3,&y\equiv0\pmod3,\\
  (2y-1)/3,&y\equiv2\pmod3.
 \end{cases}
\end{equation}
This follows immediately by separating even and odd inputs of $F$.
Every child in Theorem~\ref{thm:mnormal} lies in this image.  Represent the
$m$-coordinate $F(q)$ simply by $q$.

\begin{theorem}[Conjugated game; \proved]\label{thm:conjugate}
After the first move, the game is exactly the game on $q\in\N$ with terminal
state $0$ and moves
\[
  \boxed{A(q)=F(q)=\left\lceil\frac{3q}{2}\right\rceil,
  \qquad B(q)=R(F(q)).}
\]
Moreover, resolving every $q$-position resolves every original odd position.
\end{theorem}

\begin{proof}
Apply Theorem~\ref{thm:mnormal} to a state whose $m$-coordinate is $F(q)$.
Its child coordinates are $F(F(q))$ and $F(R(F(q)))$; after applying
$F^{-1}$ these are $A(q)$ and $B(q)$.  The terminal coordinate is
$0=F(0)$.  For an arbitrary original state with coordinate $m$, its two
children have coordinates $F(m)$ and $F(R(m))$, so they are represented by
$m$ and $R(m)$.  If every represented position is resolved, both children of
the original state are resolved, and Proposition~\ref{prop:recursion} resolves
the parent.
\end{proof}

\begin{lemma}[Strict contraction of $B$; \proved]\label{lem:Bcontract}
For every $q>0$, $B(q)<q$.
\end{lemma}

\begin{proof}
Deleting a nonempty suffix gives $R(x)\le\lfloor x/2\rfloor$ for $x>0$.
Thus, for $q\ge2$,
\[
 B(q)\le\left\lfloor\frac{F(q)}2\right\rfloor
 \le\left\lfloor\frac{3q+1}{4}\right\rfloor<q.
\]
Directly, $A(1)=2=10_2$ and $B(1)=0$.
\end{proof}

\begin{lemma}[A short descending block; \proved]\label{lem:BBA}
For every $q>0$,
\[
 B(B(A(q)))<q,
\]
where a move reaching $0$ is left at $0$.
\end{lemma}

\begin{proof}
From the proof of Lemma~\ref{lem:Bcontract},
$B(x)\le(3x+1)/4$ and $A(q)\le(3q+1)/2$.  Therefore
\[
 B(B(A(q)))\le \frac{27q+23}{32}<q
\]
for $q\ge5$.  For $q=1,2,3,4$, direct binary suffix deletion gives the
respective values $0,0,1,1$, all strictly below $q$.
\end{proof}

\begin{remark}[Why this is not a termination proof]
Lemma~\ref{lem:BBA} controls a path segment containing two specified
contracting moves.  It says nothing about arbitrarily long runs of $A$-moves,
and it does not determine which child is compatible with $\Draw$.  Therefore
it cannot by itself exclude an infinite optimal-play path.
\end{remark}

\section{The suffix operation in valuation and Gray coordinates}

\begin{lemma}[Valuation formula; \proved]\label{lem:Rvaluation}
For every $x>0$,
\begin{equation}\label{eq:Rvaluation}
 R(x)=
 \begin{cases}
 \left\lfloor x/2^{\vTwo(3x+1)}\right\rfloor,&x\text{ odd},\\[2mm]
 \left\lfloor x/2^{\vTwo(3x+2)}\right\rfloor,&x\text{ even}.
 \end{cases}
\end{equation}
\end{lemma}

\begin{proof}
Let the maximal alternating suffix have length $k$.  If $x$ is odd, its low
bits begin $1,0,1,0,\ldots$; if $x$ is even, they begin
$0,1,0,1,\ldots$.  The finite geometric sums give, modulo $2^k$,
\[
 x\equiv -3^{-1}\pmod{2^k}\quad(x\text{ odd}),
 \qquad
 x\equiv -2\cdot3^{-1}\pmod{2^k}\quad(x\text{ even}).
\]
Thus $2^k$ divides $3x+1$ in the odd case and $3x+2$ in the even case.  If a
nonempty prefix remains, the boundary bit repeats the top suffix bit; using
the same geometric sum modulo $2^{k+1}$ shows that divisibility by
$2^{k+1}$ fails.  Hence the displayed valuation equals $k$, and division by
that power of two deletes exactly the suffix.  If the suffix is the entire
binary word, the valuation may exceed $k$, but both the quotient's floor and
$R(x)$ are $0$.  This proves~\eqref{eq:Rvaluation}.
\end{proof}

\begin{lemma}[No fixed modulus; \proved]\label{lem:no-fixed-modulus}
For every $K\ge1$ there exist $x\equiv y\pmod{2^K}$ with $R(x)\ne R(y)$.
Consequently, no fixed residue modulo $2^K$ determines $R$ globally.
\end{lemma}

\begin{proof}
Fix the same alternating word, with low bit $1$, in the lowest $K$ positions
of $x$ and $y$.
For $x$, choose the bit in position $K$ to repeat the bit in position $K-1$;
for $y$, choose it to be the opposite bit and choose the next bit to repeat
it.  Add one further leading $1$ if either finite word would otherwise have a
leading zero.  Then the two integers agree in their lowest $K$ bits.  The
maximal alternating suffix of $x$ has length $K$, whereas that of $y$ has
length $K+1$.  After deletion, the low bit of $R(x)$ is the bit in position
$K$, while the low bit of $R(y)$ is the opposite bit in position $K+1$.
The retained prefixes are nonzero after the possible leading-bit addition;
their parities differ, so $R(x)\ne R(y)$.
\end{proof}

Let
\[
  G(x)=x\xor\left\lfloor\frac{x}{2}\right\rfloor
\]
be reflected binary Gray code.

\begin{theorem}[Gray deletion formula; \proved]\label{thm:gray}
For every $x\ge0$,
\[
 \boxed{G(R(x))=
 \left\lfloor
 \frac{G(x)}{2^{\vTwo(G(x)+1)+1}}
 \right\rfloor.}
\]
\end{theorem}

\begin{proof}
Write the binary digits of $x$ as $b_i$.  Away from the leading boundary, the
Gray digit in position $i$ is $b_{i+1}\xor b_i$.  If the maximal alternating
suffix has length $k$ and a nonempty prefix remains, the lowest $k-1$ Gray
digits are $1$ and the boundary Gray digit is $0$.  Thus
$\vTwo(G(x)+1)=k-1$, and shifting by $k$ leaves exactly the Gray code of the
remaining prefix.  If the entire word is alternating, $G(x)$ is an all-ones
word; the displayed shift removes the whole word (and one leading zero), so
both sides are $0$.  The case $x=0$ is immediate.
\end{proof}

\begin{proposition}[Eight-state transducer; \proved]
Read Gray words from least significant to most significant, with leading-zero
padding.  The relation
\[
  h=G\bigl(A(G^{-1}(g))\bigr)
\]
is recognized by a finite transducer with eight carry states.
\end{proposition}

\begin{proof}
Let $b_i,c_i$ be the binary digits of $q,A(q)$ and let
$g_i=b_{i+1}\xor b_i$, $h_i=c_{i+1}\xor c_i$ be their Gray digits.  Since
\[
  2A(q)=3q+b_0=q+2q+b_0,
\]
a state at column $i$ may be taken as
$\left(b_i,c_i,d_{i+1}\right)\in\{0,1\}^3$, where $d_{i+1}$ is the carry
into column $i+1$.  On reading $(g_i,h_i)$, set
\[
 b_{i+1}=b_i\xor g_i,\qquad c_{i+1}=c_i\xor h_i.
\]
The transition is legal exactly when
\[
 c_i\equiv b_{i+1}+b_i+d_{i+1}\pmod2,\qquad
 d_{i+2}=\left\lfloor\frac{b_{i+1}+b_i+d_{i+1}}2\right\rfloor.
\]
Initially $d_1=b_0$; the bits $b_0,c_0$ are guessed.  Accept after zero
padding reaches $(0,0,0)$.  Gray inversion with the leading-zero condition
makes both guesses unique.  The displayed column equations then hold if and
only if $2A(q)=3q+b_0$, proving that the accepted pairs are exactly the stated
relation.  There are $2^3=8$ states.
\end{proof}

Together with Theorem~\ref{thm:gray}, this proves that both game edges are
rational binary relations.  Finite-state representability, however, is not a
well-founded rank and does not exclude draws.

\section{Finite proof witnesses and a sound token rank}

Suppose a local outcome argument uses already certified finite proof trees.
Record only their heights in a finite multiset $M$.  The word
``already'' is essential: a finite state discovered later is not automatically
a token carried by the current configuration.

Define the numerical rank
\[
  \tau(M)=\sum_{h\in M}3^h,
\]
with multiplicity.  This is a dependency-free substitute for the natural
ordinal sum of $\omega^h$ terms when every replacement has at most two
children.

\begin{lemma}[Token replacement; \proved]\label{lem:token}
Let $M=U\uplus\{h\}$, and let
$M'=U\uplus\{h_1,\ldots,h_k\}$, where $0\le k\le2$ and every
$h_i<h$.  Then $\tau(M')<\tau(M)$.
\end{lemma}

\begin{proof}
If $k=0$ the claim is immediate.  Otherwise $h\ge1$ and
\[
 \sum_{i=1}^k3^{h_i}\le2\cdot3^{h-1}<3^h.
\]
The unchanged contribution $\tau(U)$ may be added to both sides.
\end{proof}

\begin{corollary}[Well-founded token descent; \proved]
There is no infinite sequence of finite token multisets in which every step
is a replacement of the form in Lemma~\ref{lem:token}.
\end{corollary}

\begin{proof}
The natural number $\tau(M)$ strictly decreases at every step.
\end{proof}

\begin{remark}[Provenance rule]
Lemma~\ref{lem:token} permits replacing a particular retained occurrence
$h\in M$.  It does not permit choosing an unrelated finite proof of height
$u$, proving that new heights are below $u$, and inserting them into $M$ when
$u$ was not in $M$.  For example, if $M=\varnothing$, inserting any nonempty
list increases $\tau$, regardless of comparisons with witnesses external to
$M$.
\end{remark}

\section{Three local router repairs proved explicitly}
\label{sec:local-repairs}

This section proves two draw-preserving constructions that were only named,
but not constructed, in the earlier v6.3 manuscript: the fixed-tail
loss-sibling orientation and the source-zero valuation-one exit.  It also
proves that the formally possible valuation-one branch of the shortest marked
return is empty.  None of these local lemmas repairs the separate long
high-return token-installation defect of Section~\ref{sec:audit}.

For $r\ge1$, $e\in\{0,1\}$, and odd $a>0$, put
\[
  Q_r^e(a)=2^r a-e,
  \qquad J(s)=2A(s)+1.
\]
Thus
\[
 J(s)=\begin{cases}3s+1,&s\text{ even},\\3s+2,&s\text{ odd}.
 \end{cases}
\]

\begin{lemma}[Canonical coefficient source; \proved]
\label{lem:canonical-source}
Every $x>0$ has unique coordinates
\[
 x=Q_r^e\bigl(3^kJ(s)\bigr),
 \qquad r\ge1,\quad e\in\{0,1\},\quad k,s\in\N.
\]
If $\sigma(x)=s$ denotes the resulting coefficient source, then
\[
 \sigma(x)\le \frac{x-1}{6}.
\]
\end{lemma}

\begin{proof}
The parity of $x$ determines $e$: namely $e=x\bmod2$.  The factorization
$x+e=2^r a$ with $a$ odd determines $r\ge1$ and $a$ uniquely.  Removing the
maximal power $3^k$ determines an odd integer $a_0$ not divisible by three.
Such integers are precisely $6t+1$ and $6t+5$.  The displayed formula for
$J$ gives
\[
 J(2t)=6t+1,\qquad J(2t+1)=6t+5,
\]
so $a_0=J(s)$ for a unique $s$.  Finally
\[
 x=2^r3^kJ(s)-e\ge2J(s)-1\ge6s+1,
\]
which is the asserted bound.
\end{proof}

Call $q$ \emph{exceptional} when
\[
 q\equiv1,3,12,14\pmod {16};
\]
all other $q$ are called ordinary.

\begin{lemma}[Ordinary side relation; \proved]
\label{lem:ordinary-side}
If $q$ is ordinary, then
\[
 B(A(q))\in\{A(B(q)),B(B(q))\}.
\]
Moreover, if $q$ is exceptional, then $A(q)$ is ordinary.
\end{lemma}

\begin{proof}
Put $x=A(q)$, $s=R(x)=B(q)$, and let $k$ be the length of the maximal
alternating suffix of $x$.  The condition $k\ge3$ says that the last three
bits of $x$ are $010$ or $101$.  Substitution of the sixteen residues of
$q$ into $A(q)=\lceil3q/2\rceil$ shows that this occurs exactly for
$q\equiv1,3,12,14\pmod {16}$.  Thus $k\in\{1,2\}$ when $q$ is ordinary.

If $k=1$, write $x=2s+e$, where the boundary condition for maximality says
$s\equiv e\pmod2$.  The formulas
\[
 A(2s)=3s,\qquad A(2s+1)=3s+2
\]
show that $A(x)$ is obtained from $A(s)$ by appending the bit $e$.
Deleting the new maximal alternating suffix therefore leaves $A(s)$ if the
boundary bit repeats and leaves $R(A(s))=B(s)$ if it continues.

If $k=2$, the two deleted bits are $01$ or $10$, and their high bit equals
the last bit of $s$.  Directly, $x=4s+1$ with $s$ even gives
$A(x)=4A(s)+2$, while $x=4s+2$ with $s$ odd gives
$A(x)=4A(s)+1$.  Thus $A(x)$ is obtained from $A(s)$ by appending,
respectively, $10$ or $01$.  The same boundary test leaves either $A(s)$
or $R(A(s))=B(s)$.  This proves the first assertion.  For the second,
direct substitution sends the four exceptional classes under $A$ to
$2,5,10,13\pmod {16}$, none of which is exceptional.
\end{proof}

\begin{lemma}[Exceptional child-pair attachment; \proved]
\label{lem:exceptional-pair}
Let $q$ be exceptional and put $\ell=B(q)$.  For some $m\ge1$ and
$\delta\in\{0,1\}$, the two children of $A(q)$ are exactly
\[
 Q_m^\delta(J(\ell)),\qquad Q_{m+1}^\delta(J(\ell)).
\]
Consequently, if $q$ is $\Win$, $A(q)$ is $\Draw$, and $\ell$ is $\Loss$,
then this exact adjacent pair contains a continuing $\Draw$ and remains
attached to the same losing state $\ell$.
\end{lemma}

\begin{proof}
Write $q=16t+c$, $c\in\{1,3,12,14\}$.  Direct application of $A$ and $B$
gives
\[
\begin{array}{c|c|c|c}
c&\ell&B(A(q))&A^2(q)\\ \hline
1&R(6t)&18t+1&36t+3\\
3&R(6t+1)&18t+4&36t+8\\
12&R(6t+4)&18t+13&36t+27\\
14&R(6t+5)&18t+16&36t+32.
\end{array}
\]
In the four rows put $z=6t,6t+1,6t+4,6t+5$, respectively.  If the maximal
alternating suffix of $z$ has length $m$ and leaves $\ell=R(z)>0$, direct
summation of that alternating suffix gives
\[
 3z+1=Q_m^\delta(J(\ell)),
 \qquad \delta=1-(z\bmod2).
\]
When the whole word is deleted, the same identity holds with $\ell=0$ and
$m$ one larger than the binary length.  The third column is $3z+1$ and the
fourth is twice that number plus $\delta$, proving the two coordinates.
If $A(q)$ is a draw, its children contain no loss and at least one draw by
Proposition~\ref{prop:recursion}; the final assertion follows.
\end{proof}

\begin{lemma}[Fixed-tail loss-sibling constructor; \proved]
\label{lem:fixed-tail-sibling}
Let $c>0$ be odd, $r\ge2$, and $e\in\{0,1\}$.  Put
\[
 X=Q_r^e(3c),\quad G=Q_r^e(c),\quad
 H=Q_{r+1}^e(c),\quad Y=A(G).
\]
Then $\{A(H),B(H)\}=\{X,Y\}$, in fact $A(H)=X$ and $B(H)=Y$.
Suppose
\[
 X:\Draw,\qquad H:\Win,\qquad Y:\Loss.
\]
There is an explicit draw-preserving continuation:

\begin{enumerate}
\item[(a)] if $H$ is ordinary, then $P=B(A(H))$ is simultaneously a child of
      $X$ and of $Y$; hence
      \[
      P:\Win,\qquad X\to P\text{ is a draw--win boundary},\qquad
      h(P)\le h(H)-2;
      \]
\item[(b)] if $H$ is exceptional, the two children of the actual draw $X=A(H)$
      are the exact adjacent pair of Lemma~\ref{lem:exceptional-pair} over
      the same losing state $Y=B(H)$, and that pair contains a draw.
\end{enumerate}
No transition through the winning state $G$ is asserted.
\end{lemma}

\begin{proof}
For $r\ge3$, constant-tail arithmetic gives
\[
 A(H)=Q_r^e(3c)=X,\qquad
 B(H)=Q_{r-1}^e(3c)=A(G)=Y.
\]
At $r=2$ the same second identity is the direct exponent-two/one boundary
calculation; the first identity is unchanged.

Assume first that $H$ is ordinary.  Lemma~\ref{lem:ordinary-side} gives
\[
 P=B(A(H))\in\{A(B(H)),B(B(H))\}=\{A(Y),B(Y)\}.
\]
Thus $P$ is a child of the losing state $Y$ and is winning.  It is also the
contracting child of $A(H)=X$, so it is a genuine winning child of the
continuing draw.  Since a draw cannot be the losing witness for the winning
state $H$, the stated orientation makes $Y$ its losing witness and
\[
 h(Y)=h(H)-1,\qquad h(P)\le h(Y)-1=h(H)-2.
\]
If $H$ is exceptional, apply Lemma~\ref{lem:exceptional-pair} with
$q=H$, $A(H)=X$, and $B(H)=Y$.  This proves the second alternative and
the exhaustiveness of the split.
\end{proof}

\begin{lemma}[Source-zero valuation-one exit; \proved]
\label{lem:source-zero-j1}
Let $n\ge2$, $\delta\in\{0,1\}$, and suppose
\[
 3^n+1-2\delta=2J(y)
\]
has $2$-adic valuation one.  Put
\[
 C=R(3^n-\delta),\qquad E=Q_1^{1-\delta}(J(y)).
\]
Then $y=(3^{n-1}-1)/2>0$ and exactly one of the following two coordinate
rows holds:
\[
\begin{array}{c|c|c|c}
\delta&n&C&E\\ \hline
0&\text{even}&A(y)&Q_1^1(J(y))\\
1&\text{odd}&B(y)&Q_1^0(J(y)).
\end{array}
\]
If $\{C,E\}$ contains a draw, then either $E$ itself is an explicit
exponent-one draw lift over $y$, or $C$ is a draw with
\[
 \sigma(C)<y.
\]
Thus the raw valuation-one row is a strict coefficient-source exit, not an
unconstructed loss-sibling entry.
\end{lemma}

\begin{proof}
Valuation one in $3^n+1$ forces $n$ even, while valuation one in
$3^n-1$ forces $n$ odd; this follows modulo four.  Dividing the displayed
factorization by two gives $J(y)=(3^n+1-2\delta)/2$, and the explicit
formula for $J$ yields $y=(3^{n-1}-1)/2$.  Direct substitution gives the
two signed coordinates.  In the first row the last two bits of $3^n$ are
$01$, so deletion gives $C=(3^n-1)/4=A(y)$.  In the second row appending
the final zero to $J(y)$ extends the alternating suffix of $A(y)$, giving
$C=B(y)$.

If $E$ is a draw, the first alternative is already explicit.  Otherwise a
draw in the pair must be $C$.  Lemma~\ref{lem:canonical-source} and
$A(y)\le(3y+1)/2$ give
\[
 \sigma(A(y))\le\frac{A(y)-1}{6}
 \le\frac{3y-1}{12}<y.
\]
For the other row, Lemmas~\ref{lem:Bcontract} and
\ref{lem:canonical-source} give $\sigma(B(y))<y$.  Hence both raw rows make
the claimed strict source descent.
\end{proof}

\begin{lemma}[Signed source boundary; \proved]
\label{lem:signed-source-boundary}
Let $s>0$, $e\in\{0,1\}$, and factor
\[
 3J(s)+1-2e=2^vJ(t).
\]
Then
\[
 t=\begin{cases}
 A(s),&v=1,\\
 B(s),&v\ge2.
 \end{cases}
\]
\end{lemma}

\begin{proof}
The odd integer $J(s)$ is the original game state represented by $s$.
Theorem~\ref{thm:conjugate} says that the sources of its two signed children
are exactly $A(s)$ and $B(s)$.  Of the consecutive even integers
$3J(s)-1$ and $3J(s)+1$, one has valuation one and the other has valuation at
least two.  The valuation-one odd quotient is greater than $J(s)$, whereas
the valuation-at-least-two quotient is smaller than $J(s)$.  Since $J$ is
strictly increasing, $A(s)>s>B(s)$ by Lemma~\ref{lem:Bcontract}; hence the
first quotient has source $A(s)$ and the second has source $B(s)$.
\end{proof}

\begin{lemma}[The shortest marked return is strict; \proved]
\label{lem:shortest-marked-return}
Let $C>0$ be odd, let $g\in\{0,1\}$, and put
\[
 b=Q_1^g(C),\qquad q=B(b),\qquad y=A(b).
\]
Suppose that $b,q$ are $\Win$ and $y$ is $\Loss$.  Factor
\[
 9J(b)+1-2g=2^jJ(w).
\]
Then $j\ge2$ and
\[
 w=\begin{cases}
 A(y),&j=2,\\
 B(y),&j\ge3.
 \end{cases}
\]
Consequently $w$ is an actual winning child of the retained losing state
$y$, and
\[
 h(w)\le h(y)-1.
\]
In particular this marked $D=1$ row is a legal strict token replacement; it
has no valuation-one lift subcase.
\end{lemma}

\begin{proof}
Because $b=2C-g$ has parity $g$, the formula for $J$ gives
\[
 J(b)=6C+1-2g.
\]
Hence
\[
 9J(b)+1-2g=2(27C+5-10g).
\]
The integer in parentheses is even for both values of $g$, since $C$ is odd.
Thus $j\ge2$.

It remains to identify the returned source.  First note that
\[
 y=A(b)=3C-g,
\]
and therefore
\[
 J(y)=\begin{cases}
 9C+2,&g=0,\\
 9C-2,&g=1.
 \end{cases}
\]
The remaining factorization is precisely the signed source boundary at $y$:
\[
 3J(y)+1-2(1-g)=27C+5-10g=2^{j-1}J(w).
\]
If $j=2$, its valuation is one and
Lemma~\ref{lem:signed-source-boundary} gives $w=A(y)$.  If $j\ge3$, its
valuation is at least two and the same lemma gives $w=B(y)$.

In both cases $w$ is a child of the retained $\Loss$ state $y$.
Proposition~\ref{prop:recursion} makes it $\Win$, and adjoining this child to
the finite losing proof for $y$ gives $h(w)\le h(y)-1$.
\end{proof}

\section{The abstract macro theorem}

This section states exactly what a global certificate would have to prove.
It deliberately separates a sound order-theoretic kernel from the semantic
connection to the game.

Let a macro configuration be
\[
  C=(s,M,a,N,c),
\]
where $s,a\in\N$ are outer and inner source parameters, $M,N$ are finite
multisets of retained proof heights, and $c$ belongs to a finite control set.
Assume the equal-rank control graph is acyclic, and choose a topological rank
$d(c)\in\N$ that decreases on every control edge.  Put
\[
 \rho(C)=\bigl(s,\tau(M),a,\tau(N),d(c)\bigr)\in\N^5
\]
with lexicographic order.

\begin{lemma}[Macro well-foundedness; \proved]\label{lem:macrowf}
Any relation $C'\prec C$ satisfying $\rho(C')<_{\rm lex}\rho(C)$ is
well-founded.
\end{lemma}

\begin{proof}
The usual order on $\N$ is well-founded, and a finite lexicographic product of
well-founded orders is well-founded.  Concretely, an infinite descending
sequence would first strictly decrease the first coordinate only finitely
often.  On a tail it is constant; repeat this argument successively for the
remaining four coordinates.  The resulting tail would have all coordinates
constant, contradicting strict lexicographic decrease.
\end{proof}

\begin{definition}[Complete draw-macro refinement]\label{def:refinement}
A complete refinement consists of a predicate $\mathsf{Bad}(C)$ and two
constructions:
\begin{enumerate}
\item[(R1)] for every conjugated draw $q$, a configuration $C$ with
      $\mathsf{Bad}(C)$;
\item[(R2)] for every $C$ with $\mathsf{Bad}(C)$, a configuration $C'$ with
      $\mathsf{Bad}(C')$ and $C'\prec C$.
\end{enumerate}
Here every macro edge used in (R2) must be a finite legal game construction,
must follow an actual $\Draw$ continuation supplied by
Proposition~\ref{prop:recursion}, and must preserve the identities and
provenance of all tokens unless Lemma~\ref{lem:token} certifies their strict
replacement.
\end{definition}

\begin{theorem}[Conditional no-draw theorem; \proved]\label{thm:conditional}
If a complete draw-macro refinement in the sense of
Definition~\ref{def:refinement} exists, then the conjugated game has no draw
positions.  Consequently, optimal play in the original game terminates at
$1$ from every odd starting value.
\end{theorem}

\begin{proof}
If a draw existed, (R1) would make the set of bad configurations nonempty.
By Lemma~\ref{lem:macrowf}, every nonempty subset has a $\prec$-minimal
element $C$.  But (R2) supplies a bad $C'\prec C$, a contradiction.  Hence
there are no conjugated draws.  Theorem~\ref{thm:conjugate} resolves every
original position, and Proposition~\ref{prop:strategy} gives finite
termination under optimal play.
\end{proof}

\subsection{What the routing certificate does and does not prove}

The checked inventory has 15 control states, 45 declared transitions, 12
declared total guard partitions, 29 rank-strict transitions, and 16
equal-rank transitions forming a directed acyclic graph.  These facts prove
Lemma~\ref{lem:macrowf} for the \emph{declared} macro relation.  To instantiate
Theorem~\ref{thm:conditional}, four additional obligations must be proved:

\begin{enumerate}
\item[(H1)] \textbf{Arithmetic validity.}  Every declared row has the asserted
      exact child, source, exponent, and suffix identities.
\item[(H2)] \textbf{Semantic coverage.}  The guards partition every game situation
      that can arise from a draw; there is no missing entry or reset case.
\item[(H3)] \textbf{Outcome and token compatibility.}  The selected route follows
      an actual draw child, and every rank comparison concerns the actual
      retained token occurrences, not newly introduced finite witnesses.
\item[(H4)] \textbf{Finite productivity.}  Each macro invocation reaches a declared
      successor after finitely many legal constructions, and a reset occurs
      only after a preceding rank component has strictly decreased.
\end{enumerate}

The JSON checker validates declarations about the abstract inventory.  It
cannot establish (H2)--(H4) merely by confirming that the declared guards are
syntactically disjoint and exhaustive over their own declared parameter
domains.  Likewise, the Lean structure called
\texttt{DrawMacroRefinement} makes (R1)--(R2) explicit hypotheses; no concrete
value of that structure is currently constructed.

\section{Hostile audit of the proposed global bridge}\label{sec:audit}

This section records proof defects, not merely requests for more exposition.
Defect A applies to the current human assembly.  Defect B was a defect of the
earlier v6.3 manuscript; it is repaired explicitly by
Lemma~\ref{lem:fixed-tail-sibling}.  The source-zero valuation-one omission
from that manuscript is repaired by Lemma~\ref{lem:source-zero-j1}.  The
current proof still cannot instantiate Theorem~\ref{thm:conditional} until
Defect A and the explicit semantic-refinement boundary are closed.

\subsection{Defect A: the high-return witnesses are not installed tokens}

In the high-return module, the argument constructs finite winning states
$u,c$ and lower descendants $p,b$, with height inequalities of the schematic
form
\[
  h(p)<h(u),\qquad h(b)<h(c).
\]
The next assembly step treats this as a strict replacement
\[
  \{h(u),h(c)\}\longmapsto\{h(p),h(b)\}
\]
inside the current token multiset.  However, the local lemma proves only that
$u$ and $c$ are finite winning states; it does not prove that occurrences of
$h(u)$ and $h(c)$ were already retained in the current outer or inner
multiset.

This is a logical failure of the premise of Lemma~\ref{lem:token}.  If the
current multiset is empty, for example, installing $h(p)$ and $h(b)$ makes
$\tau$ positive rather than smaller.  Comparisons with external witnesses do
not repair the problem.  A valid repair must either:

\begin{enumerate}
\item[(a)] trace $u$ and $c$ as exact retained token occurrences from the entry
      configuration to the high return; or
\item[(b)] define a different rank that pays for their installation before the
      return and prove strict descent for every reset.
\end{enumerate}

No such trace or prepaid rank is currently supplied.  Therefore the claimed
strict high-return edge is \openstatus.

\subsection{Repaired defect B (v6.3 only): the loss-sibling orientation}

One entry pattern in the v6.3 normalizer has the outcome fingerprint
\[
  X:\Draw,\qquad H:\Win,\qquad Y:\Loss,\qquad G:\Win.
\]
The available local diamond keeps the draw at $X$.  The assembly nevertheless
routes the configuration through a ``loss-sibling entry'' based at $G$.
Because $G$ is winning rather than drawing, the v6.3 transition was not an
outcome-compatible continuation of $X$.  Lemma~\ref{lem:fixed-tail-sibling}
now supplies the missing constructor without moving the draw to $G$.  In the
ordinary row, the actual draw $X$ has a winning child two proof levels below
$H$; in the exceptional row, the two actual children of $X$ form an exact
adjacent frame over the same losing state $Y$.  Thus this historical defect
is closed in the present manuscript.

\subsection{Additional obligations exposed by the audit}

Two further points in v6.3 required explicit case proofs in addition to
Defect B:

\begin{itemize}
\item the $D=2$ loss-anchored lift/factor exits must preserve or lower the
      exact incoming token in every signed and raw exit, not merely produce
      some finite witness; this local obligation is now closed by the exact
      split $B(w)\in\{A^2(q),B(A(q))\}$ and the ordinary/exceptional
      common-child identities recorded in the repository note
      \emph{D2 loss-anchored obligation}, conditional only on the incoming
      occurrence of $q$ being legally installed;
\item the source-zero, valuation-one branch required an explicit route;
      Lemma~\ref{lem:source-zero-j1} now proves that the signed exit is an
      exact exponent-one lift and the raw exit is a strict coefficient-source
      descent, so no new loss-sibling constructor is needed.
\end{itemize}

None of these observations disproves the desired game theorem.  Defect A shows
that the current token installation and lifecycle have not yet been justified.
The loss-sibling, source-zero, shortest marked $D=1$, and local $D=2$ rows
are now closed.  The long high-return carry/no-reseed theorem remains the
explicit global obligation.
Independently, the current Lean development leaves the complete draw-to-macro
map as an explicit structure parameter.  Thus the last implication needed
for an unconditional proof is unavailable.  The exact acceptance conditions
for Defect A are recorded in the repository note
\emph{High-return provenance obligation}.

\section{Computational verification and its exact scope}\label{sec:computation}

All commands below were run on 12 August 2026 using Python 3.12.13.  They are
reproducible from the repository snapshot audited for this manuscript.

\begin{verbatim}
python -m unittest discover -s tests -v
python scripts/verify_claims.py --limit 100000
python scripts/verify_global_certificate.py
python audit.py --limit 100000
\end{verbatim}

The results were:

\begin{itemize}
\item 112 unit and regression tests passed;
\item the $m$-space normal form was checked for $1\le m\le100{,}000$;
\item $B(q)<q$ was checked for $1\le q\le100{,}000$;
\item descent blocks were checked for odd $3\le n\le200{,}001$;
\item the global inventory checker accepted 15 states and 45 transitions,
      reporting the status \texttt{CONDITIONAL\_MACHINE\_CHECK} and four
      unverified local refinement/outcome obligations;
\item an independently verified finite proof DAG certified $q=100$ as
      winning, using 59 nodes and maximum finite rank 44,448.
\end{itemize}

The previous manuscript v6.3 was also tested by
\begin{verbatim}
python verify_structure_v6_3.py
python verify_v6_3.py --limit 250000 \
  --random-samples 5000 --random-bits 256
\end{verbatim}
Both scripts passed.  The second script explicitly describes its scope as
``arithmetic/control regressions only; human proof carries outcomes/ranks.''
This is consistent with Defects A and B: neither is a false integer identity
detectable by a finite arithmetic test.

These results are \verified, not theorems about all integers.  In particular,
an unknown node outside a finite retrograde window is not labelled winning or
losing, and no finite cutoff is treated as a boundary condition.

\section{A precise target for future work}

The open problem has now been reduced to a concrete construction rather than
an informal appeal to a complicated case analysis.

\begin{proposition}[Sufficient remaining deliverable; \proved]
It is sufficient to provide, in Lean or on paper, explicit functions
implementing (R1)--(R2) of Definition~\ref{def:refinement}, together with
proofs of (H1)--(H4).  No additional global game-theoretic argument would then
be needed.
\end{proposition}

\begin{proof}
Such functions instantiate the hypotheses of
Theorem~\ref{thm:conditional}, whose conclusion is exactly the desired
no-draw statement and, by conjugacy, optimal termination in the original
game.
\end{proof}

For a hostile-audit-ready completion, the refinement should expose at least
the following data on every transition:

\begin{enumerate}
\item the current legal game states and their certified outcomes;
\item the actual draw child selected by Proposition~\ref{prop:recursion};
\item the identities of retained token occurrences before and after the
      transition;
\item the exact arithmetic guard and proof that all alternatives are covered;
\item the first coordinate of $\rho$ that strictly decreases, or the verified
      equal-rank control edge;
\item a finite bound or structural recursion proving that the macro returns.
\end{enumerate}

This format would make the two defects in Section~\ref{sec:audit} impossible
to hide: an absent parent token cannot be named in item 3, and a winning state
cannot be supplied as the draw child in item 2.

\section{Conclusion}

The exact binary conjugacy and the finite-outcome semantics isolate the real
difficulty of the two-player $3n\pm1$ game.  The contracting branch, the
Gray-code transducer, finite arithmetic checks, and a well-founded abstract
macro rank are all useful and rigorously established.  They do not by
themselves show that every draw continuation enters that rank system.

The current global proof attempt fails at a specific semantic link: it uses
high-return witnesses that have not been installed as retained tokens.  The
earlier v6.3 attempt additionally routed one loss-sibling orientation through
a winning state rather than the actual draw.  Accordingly, the unconditional
theorem remains \openstatus.  Theorem~\ref{thm:conditional} states a complete and
machine-friendly criterion whose future instantiation would finish the
problem without any further appeal to informal global induction.

\appendix
\section{Adversarial-audit checklist}

Before a future version claims the global target as \proved, each answer
below must be supported by a proof object, a theorem
reference, or an explicit finite case table with symbolic coverage.

\begin{enumerate}
\item[\textbf{A1.}] Is every use of ``draw has a draw child'' attached to the actual child
      selected in the next macro state?
\item[\textbf{A2.}] Is every token decreased by identity, not only by comparing numerical
      heights of unrelated finite states?
\item[\textbf{A3.}] When a seed is installed, which older rank coordinate has already
      decreased and pays for the reset?
\item[\textbf{A4.}] Are source and coefficient kept distinct under every $3^k$ factorization?
\item[\textbf{A5.}] Are all suffix lengths handled symbolically, rather than by a fixed
      modulus or search depth?
\item[\textbf{A6.}] Does every loss-sibling orientation have an explicit constructor?
\item[\textbf{A7.}] Are the short exponents $1,2,3$ proved separately wherever the generic
      descent uses $D\ge3$ or $D\ge4$?
\item[\textbf{A8.}] Does every macro return after finitely many steps, including exceptional
      and source-zero branches?
\item[\textbf{A9.}] Are equal-rank edges acyclic after preserving the actual full token
      multisets, not merely their numerical ranks?
\item[\textbf{A10.}] Has the concrete refinement, rather than only the abstract rank theorem,
      been checked by the proof assistant?
\end{enumerate}

\begin{thebibliography}{9}

\bibitem{Guy1986}
R.~K. Guy,
\emph{John Isbell's Game of Beanstalk and John Conway's Game of
Beans-Don't-Talk},
Mathematics Magazine \textbf{59} (1986), 259--269.
\url{https://doi.org/10.1080/0025570X.1986.11977258}.

\bibitem{Guy1996}
R.~K. Guy,
\emph{Unsolved Problems in Combinatorial Games}, Problem 42,
in Games of No Chance, MSRI Publications 29, Cambridge University Press,
1996.

\bibitem{AHZ2024}
I.~Alth\"ofer, M.~Hartisch, and T.~Zipproth,
\emph{Analysis of a Collatz Game and Other Variants of the $3n+1$ Problem},
Advances in Computer Games (2024), 123--132.
\url{https://doi.org/10.1007/978-3-031-54968-7_11}.

\bibitem{Repository}
G.~Pochuev,
\emph{Optimal $3n\mathbin{\pm}1$ Game Repository}, source code and proof
notes, 2026.  \url{https://github.com/Grisha-Pochuev/3n-plus-minus-1-game}.

\end{thebibliography}

\end{document}
